\documentclass[11pt,a4paper]{article}
\usepackage[margin=24mm]{geometry}
\usepackage{amsmath,amssymb,amsthm,booktabs,hyperref}
\newtheorem{theorem}{Theorem}\newtheorem{proposition}[theorem]{Proposition}
\newcommand{\gen}[1]{\langle #1\rangle}
\title{Explicit coset examples and their full triality centralizers}
\author{Human Compression project: recovery of L122}\date{}
\begin{document}\maketitle
The finite examples of Leemans and Stokes~\cite{LS} are genuine mathematical
claims, including their full automorphism groups. The original does not give
its subgroup markings or software. We supply complete explicit witnesses
for every listed parameter. We do not identify these witnesses with the
unprovided historical choices. The complete companion \emph{Coset trialities,
chamber recognition and reduced relations} supplies the construction,
flag-completion and full-automorphism arguments used below.

\section{Exact finite-field markings}
For each field below, $z$ is the residue of $X$ in the indicated polynomial
quotient. These are the primitive Conway generators used by GAP; the
polynomial coefficients also specify the fields independently of that name.
\[
\begin{array}{c|l}
q&\text{defining polynomial}\\\hline
8&X^3+X+1\\
27&X^3+2X+1\\
64&X^6+X^4+X^3+X+1\\
125&X^3+3X+3\\
512&X^9+X^4+1
\end{array}
\]
For $q=8,27,64,125,512$, use the natural projective action of
$\operatorname{SL}_2(q)$ on the $q+1$ lines represented by $(1,x)$ and
$(0,1)$. Its image is $G=\operatorname{PSL}_2(q)$. The following
matrices of determinant one specify the first projective involution $r_0$:
\[
\begin{array}{c|c}
q&R_0\\\hline
8&\begin{pmatrix}1&0\\1&1\end{pmatrix}\\[2mm]
27&\begin{pmatrix}z^2&1\\z^5&z^{15}\end{pmatrix}\\[2mm]
64&\begin{pmatrix}z^{60}&z^{11}\\z^9&z^{60}\end{pmatrix}\\[2mm]
125&\begin{pmatrix}z&1\\z^9&z^{63}\end{pmatrix}\\[2mm]
512&\begin{pmatrix}z^{507}&z^9\\z&z^{507}\end{pmatrix}
\end{array}
\]
For $q=27,64,125,512$, put respectively $p=3,4,5,8$ and let $\sigma$
be the field automorphism raising every entry to its $p$th power.
Set $r_i=\sigma^i(r_0)$ for $i=0,1,2$. For $q=8$ use instead conjugation
by the projectivity represented by
\[
 T=\begin{pmatrix}z^5&z^5\\z&z^4\end{pmatrix},
\]
which has projective order three. This preserves the source's explicit
inner/outer distinction. Set $H_i=\gen{r_j,r_k}$ for the complementary pair.

For the Suzuki example use the standard four-dimensional matrix group
$\operatorname{Sz}(8)$, acting on its 65-point ovoid, with the same generator
$z\in\mathbb F_8$. Here take
\[
 R_0=\begin{pmatrix}
z^6&z^3&z^4&1\\ z^6&z^6&z^3&z^4\\
z^5&z&z^6&z^3\\ z^4&z^5&z^6&z^6
\end{pmatrix}
\]
and let $\sigma$ square every entry. Again use its three cyclic images and
the complementary pair subgroups. Standard Suzuki coordinates and the
faithful ovoid identification are retained classical inputs, already used
in the earlier Suzuki work; the exact standard matrices and ovoid are
those of Bray and B\"a\"arnhielm~\cite{BB}, equations (1)--(5); no group is identified merely by its order.

\section{The whole finite certificate and geometric conclusions}
The required file \path{programs/verify_examples.g}, together with the
whole marking data \path{data/seeds.g}, checks the six cases. For the
projective linear groups it constructs every projective point and its
matrix action. For the Suzuki group it constructs the full 65-point orbit
of $(1,0,0,0)$. In every case the permutation order is verified on a fresh
group formed solely from copies of the permutation generators, with no
cached size, chain, ambient parent or supplied order bound. Stabilizer-chain
verification uses \texttt{random=1000}; internal random construction is not
excluded. The matrix group models, exact finite-field and subgroup operations
are explicit GAP trust dependencies.

The program checks every seed's membership and involution order, full
generation, order-three automorphism, cyclic transport of all three
parabolics, all 64 subset intersections, and each of the three product
intersections $H_i\cap H_jH_k$. The latter have exactly four elements.
It obtains the following full results; $m=|r_0r_1|$, and $v$ is the number
of objects of each type.
\begin{center}
\begin{tabular}{lrrrrr}
\toprule
$G$&$|G|$&$|H_i|$&$m$&$v$&$|G^\sigma|$\\\midrule
$\operatorname{Sz}(8)$&29120&14&7&2080&20\\
$\operatorname{PSL}_2(8)$&504&14&7&36&9\\
$\operatorname{PSL}_2(27)$&9828&28&14&351&12\\
$\operatorname{PSL}_2(64)$&262080&130&65&2016&60\\
$\operatorname{PSL}_2(125)$&976500&62&31&15750&60\\
$\operatorname{PSL}_2(512)$&134217216&1026&513&130816&504\\\bottomrule
\end{tabular}
\end{center}
These values and the verified equalities are computational proof premises;
there is no reliance on an unprovided historical Magma run. The subsequent
passage to every flag and every automorphism is written, not sampled.

\begin{theorem}\label{examples}
For all six markings the coset system is a thin, residually connected,
flag-transitive rank-three geometry with a triality induced by $\sigma$.
Its full type-preserving automorphism group is $G$. Its full reduced oriented
graph automorphism group is $G^\sigma$.
\end{theorem}
\begin{proof}
The certificates give exactly the hypotheses of the companion's chamber
recognition theorem, including trivial chamber stabilizer and the complete
four-element product identity. That theorem proves all the geometric
properties and that no additional typed automorphism exists. The cyclic
parabolic transport and $\sigma^3=1$ give an actual triality. The companion's
reconstruction and centralizer theorem then proves the last assertion.
No inference from the order of a discovered automorphism subgroup to the
order of the full graph group is made.
\end{proof}

The outer qualification is also checked. On each natural point set the
program constructs the field permutation, verifies its conjugation on all
group generators, checks that it is outside $G$, and verifies that the
point stabilizer is transitive on the remaining points. A doubly transitive
permutation group on at least three points has trivial centralizer: a
centralizing permutation sending $x$ to a different $y$ would force the
point stabilizer at $x$ to fix $y$, contrary to that transitivity. A
centralizing permutation fixing $x$ fixes every point. Thus a field
permutation outside $G$ cannot induce an inner automorphism. This proves
the four projective-linear and the Suzuki outer assertions.

For $q=p^3$ in the four outer projective linear cases, the fixed group is
precisely $\operatorname{PSL}_2(p)$. Indeed if a projectivity represented by
$A\in\operatorname{SL}_2(q)$ is fixed, then $A^{(p)}=\epsilon A$, where
$\epsilon=\pm1$ in odd characteristic and $\epsilon=1$ in characteristic
two. Three iterations give $\epsilon^3=1$, so in both cases $\epsilon=1$.
Thus every entry belongs to $\mathbb F_p$, and the converse is immediate.
The scalar-center quotient gives exactly the asserted group. The program
also constructs the embedded $\operatorname{SL}_2(p)$ image, verifies its
order and checks that each of its generators is fixed.
For $q=8$ the centralizer of the specified inner element is cyclic of order
nine; the program checks its actual subgroup and cyclicity. This explains
the source's exception to the subfield pattern.

For the Suzuki case, the field-fixed group is $\operatorname{Sz}(2)$ of
order 20. This can be seen from the unique standard Bruhat parameters:
matrices fixed by squaring have parameters in $\mathbb F_2$, giving
$2^2+2^4=20$ elements. This is the same retained matrix/Bruhat identification
used in the earlier Suzuki field-intersection proof. The certificate independently checks all 29120 permutation elements, finding
exactly 20 fixed ones. It also constructs the binary-coordinate subgroup, checks its order 20, its normal
cyclic subgroup of order 5, and a cyclic complement of order 4 acting faithfully.
It is therefore $C_5:C_4$, as required, rather than a dihedral group of order 20.

\section{The absolute subgeometry in the Suzuki case}
An object $H_0g$ is absolute precisely when
\[
 g\sigma(g)^{-1}\in H_0H_1.
\]
This is just the right-coset intersection criterion for $H_0g$ and
$H_1\sigma(g)$. The program tests every one of the 2080 cosets and obtains
exactly 20 absolute objects. The order 20 fixed subgroup has one orbit on them,
necessarily regular. By reconstruction each produces a reduced loop and a
triple of pairwise incident objects in the original geometry.
The induced system on all 60 of these objects is checked in full: it has
60 edges and 20 chambers, and every edge extends. More precisely it is a
disjoint union of 20 chambers. Thus it is a subgeometry with the stated
fixed-subgroup action; connectedness of this subgeometry was not claimed by
the source and is not silently added. The full incidence graph has 6240
vertices and the reduced oriented graph 2080 vertices, as in~\cite{LS}.

\section{The Mathieu example and why the full group matters}
Retain the classical degree-twelve, three-transitive action of $M_{11}$
on the cosets of a $\operatorname{PSL}_2(11)$ subgroup, and the classical
fact $\operatorname{Aut}(M_{11})=M_{11}$. These named-group facts are separate
from the following elementary incidence calculation.

\begin{proposition}\label{mathieu}
Let $G=M_{11}$ in that action on $\Omega$, $|\Omega|=12$, and take
$H_i$ to stabilize three distinct points. The resulting geometry consists
of three typed copies of $\Omega$, with distinct labels incident. It is
firm, residually connected and flag-transitive, has full typed group
$S_{12}$ and full type-permuting group $S_{12}\times S_3$, and admits a
triality induced by an inner automorphism of $G$.
\end{proposition}
\begin{proof}
Identify a coset of a point stabilizer with the image of its point.
Two cosets meet exactly when one element of $G$ realizes both prescribed
images. Two-transitivity says this happens exactly for distinct target
labels. Three-transitivity gives transitivity on chambers, which are
precisely triples of distinct labels. Every flag extends and every panel
has ten completions, proving firmness.
A vertex residue is the bipartite graph on two copies of the other eleven
labels, with all unequal pairs joined. This graph is connected: any two
vertices in the same part have a common neighbor, and the opposite-part
case follows by a path of length at most three. The full graph is connected
as well. These are all residues of rank at least two.

A typed automorphism is a triple of permutations of $\Omega$. Nonincidence
between two different types is exactly equality of labels, so the three
permutations must coincide. Conversely every simultaneous permutation
preserves incidence. Arbitrary permutations of the three types also
preserve incidence and commute with the simultaneous label permutations.
This proves the two full group assertions, including the absence of any
additional correlations.

There is an element of order three in $M_{11}$, and its faithful action has
a three-cycle. By three-transitivity, conjugate that cycle to the three
chosen points. Conjugation by this element cyclically permutes their
stabilizers and has order three. The coset construction therefore gives
the required inner triality. The larger $S_{12}$ typed group shows why
one cannot replace the full-group proof of Theorem~\ref{examples} by merely
exhibiting the original coset action.
\end{proof}

\section{Source contract and remaining qualification}
All three finite example sections of~\cite{LS} are now supplied with
explicit models and proofs of their full-group claims. The historical
software mentioned there is Magma, Sage and GAP; no source executable
listing is supplied. Our complete required certificate and data are
separate from any future manuscript-length comparison. No software cost
is declared zero, and no factor or full formalization is claimed.

The companion corrects the later higher-order normalization assertions with
explicit counterexamples and gives the lossless order-$r$ voltage convention.
The stronger literal reading of Lemma~1 without an order condition remains
separate: cycling parabolics proves an induced correlation with the specified
type cycle, but the argument alone does not show order three. Our examples
satisfy the necessary order condition. No complete-source admission is
claimed while the general bare existential reading is unresolved. The source's
proposed future exploration of other $r$-alities is not treated as a solved
question, and no historical novelty is inferred from these repairs.

\begin{thebibliography}{9}
\bibitem{LS} D. Leemans and K. Stokes, \emph{Coset geometries with trialities
and their reduced incidence graphs}, Acta Math. Univ. Comenianae
\textbf{88}(3) (2019), 911--916.
\bibitem{Suzuki} M. Suzuki, \emph{A new type of simple groups of finite order},
Proc. Natl. Acad. Sci. USA \textbf{46} (1960), 868--870.
\bibitem{Wilson} R. A. Wilson, \emph{On the simple groups of Suzuki and Ree},
Proc. Lond. Math. Soc. \textbf{107} (2013), 680--712.
\bibitem{BB} J. N. Bray and H. B\"a\"arnhielm, \emph{A new method for recognising Suzuki groups}, arXiv:1706.05697v1 (2017), equations (1)--(5). Auxiliary classical model, not an additional original.
\bibitem{GAP} The GAP Group, \emph{GAP}, version 4.15.1. Required exact
permutation, finite-field and subgroup operations; named model identities
and full runtime costs remain explicit external inputs.
\end{thebibliography}
\end{document}
